\begin{figure}[h]
\centering
\begin{tikzpicture}[
    scale=1,
    box/.style={draw, rounded corners=4pt, minimum width=2.8cm, 
                minimum height=1cm, align=center, font=\small},
    solidbox/.style={box, fill=gray!20},
    fluidbox/.style={box, fill=blue!10},
    homogbox/.style={box, fill=orange!15},
    arrow/.style={-{Stealth}, thick},
    label/.style={font=\footnotesize, align=center}
]

%--- Bloc 1 : Échelle du pore (image PDF en fond) ---
\begin{scope}[xshift=0cm]
    % Image de fond à la place du rectangle gris
    \node[inner sep=0pt, anchor=south west] at (0,0) 
        {\includegraphics[width=2.8cm,height=3.5cm]{figures/fig2_3D}};
    % Titre
    \node[below, label] at (1.4, -0.1) {3D hétérogène\\(échelle du pore $\ell_p$)};
\end{scope}

%--- Flèche + label homogénéisation ---
\begin{scope}[xshift=3.2cm]
    \draw[arrow, double, double distance=2pt] (0, 1.75) -- (1.4, 1.75);
    \node[above, font=\footnotesize, align=center] at (0.7, 1.9) 
        {Prise de};
    \node[below, font=\footnotesize, align=center] at (0.7, 1.6) 
        {moyenne};
\end{scope}


%--- Bloc 2 : VER (Volume Élémentaire Représentatif) ---
\begin{scope}[xshift=4.8cm]
    % Fond homogène
    \fill[orange!10, rounded corners=3pt] (0,0) rectangle (2.8,3.5);
    \draw[orange!60!black, thick, rounded corners=3pt] (0,0) rectangle (2.8,3.5);
    % Propriétés effectives
    \node[font=\scriptsize, align=center, orange!70!black] at (1.4, 2.6) 
        {$\varepsilon_l$, $a_v$};
    \node[font=\scriptsize, align=center] at (1.4, 2.0) 
        {$D^{eff}_i$, $\sigma^{eff}$};
    \node[font=\scriptsize, align=center] at (1.4, 1.4) 
        {$\langle C_i \rangle^l$, $\langle \varphi_l \rangle^l$};
    \node[font=\scriptsize, align=center] at (1.4, 0.8) 
        {$\langle \varphi_s \rangle^s$};
    % Titre
    \node[below, label] at (1.4, -0.1) {Milieu homogène\\équivalent};
\end{scope}

%--- Flèche + label réduction 1D ---
\begin{scope}[xshift=8.0cm]
    \draw[arrow, double, double distance=2pt] (0, 1.75) -- (1.4, 1.75);
    \node[above, font=\footnotesize, align=center] at (0.7, 1.9) 
        {Réduction};
    \node[below, font=\footnotesize, align=center] at (0.7, 1.6) 
        {1D};
\end{scope}

%--- Bloc 3 : Modèle 1D ---
\begin{scope}[xshift=9.6cm]
    % Fond 1D
    \fill[green!10, rounded corners=3pt] (0.9,0) rectangle (1.9,3.5);
    \draw[green!60!black, thick, rounded corners=3pt] (0.9,0) rectangle (1.9,3.5);
    % Axe z
    \draw[-{Stealth}, gray!60] (3,3.8) -- (3,-0.3);
    \node[right, font=\scriptsize] at (3, -0.3) {$z$};
    % Graduation en haut (paroi, z = -l_e)
    \draw[gray!60] (2.8,3.4) -- (3.2,3.4);
    \node[left, font=\tiny, gray!70!black] at (2.9,3.4) {$-\ell_e$};
    % Graduation en bas (bulk, z = 0)
    \draw[gray!60] (2.8,0) -- (3.2,0);
    \node[left, font=\tiny, gray!70!black] at (2.9,0) {$0$};
    % Conditions limites
    \node[left, font=\tiny, align=right] at (0.9, 3.4) 
        {Métal};
    \node[left, font=\tiny, align=right] at (0.9, 0.2) 
        {Phase homogène};
    % Titre
    \node[below, label] at (1.4, -0.1) {Modèle 1D\\(échelle $L$)};
\end{scope}

%--- Accolade échelles ---
\draw[decorate, decoration={brace, amplitude=6pt, mirror}, gray!60]
    (0,-1.2) -- (7.7,-1.2) 
    node[midway, below=8pt, font=\footnotesize, gray!70!black] 
    {$\ell_p \ll L$};

\end{tikzpicture}
\caption{Démarche d'homogénéisation : passage du modèle 3D hétérogène 
à l'échelle du pore au modèle 1D homogénéisé à l'échelle de l'électrode.}
\label{fig:homogeneisation}
\end{figure}



